\chapter[Preprocessing and Memory Engineering]{Preprocessing, Normalization, and Memory Engineering}
\label{ch:preprocessing}

Preprocessing decisions look clerical and are architectural. The scaler you
choose interacts with drift; the precision you store in determines what
experiments are affordable; the row order on disk decides whether a sequence
model trains in minutes or hours. This chapter covers the unglamorous layer
where several of our campaign's real bottlenecks lived.

\section{Know your panel's pathologies first}

Before transforming anything, profile the raw table. Ours had four
pathologies, each demanding a decision:

\begin{itemize}
\item \textbf{Ragged early history.} The number of timestamps per day drifts
  in early data and stabilizes at 968 only from date 677. Sequence models
  that batch a day as a fixed-shape tensor need the invariant; the reference
  solution (and ours) simply drops dates below 700. Losing 40\% of rows
  sounds expensive; under drift (\S\ref{sec:drift}) the old rows were
  worth little anyway --- our window sweeps later confirmed the discard was
  free.
\item \textbf{Missing symbols.} Not every symbol trades every day. Every
  reshape to (symbols $\times$ time) must handle absent (symbol, day)
  blocks explicitly --- we keep only complete blocks per day and let the
  cross-sectional features carry the survivors' information.
\item \textbf{Nulls with structure.} Feature nullness is not random: some
  features are null for entire symbols or early eras. Our profiler tracks
  a null-rate column per feature; features crossing 50\% null in a window
  are excluded from that window's statistics rather than imputed into
  fiction. Where imputation is needed for tensor models we use 0 (the
  post-standardization mean), \emph{after} the standardizer --- imputing
  before fitting statistics contaminates them.
\item \textbf{Heavy tails.} Financial features have outliers that are real
  but rare. Left unclipped they dominate variance estimates and gradients
  (Chapter~\ref{ch:metric}); clipped too aggressively they lose tail
  information that trees, in particular, can use.
\end{itemize}

\section{Normalization: which invariance do you want?}

Every normalization buys an invariance and pays with a piece of information.
Choose deliberately, per feature family, not globally:

\begin{center}
\begin{tabularx}{\textwidth}{@{}lXX@{}}
\toprule
transform & buys & pays \\
\midrule
global standardize (frozen) & comparable scales for NN optimization &
none, if statistics are trailing-legal \\
per-symbol rolling mean removal & removes slow symbol-specific level;
adapts to drift & destroys the level (fine if level is stale; fatal if
level is the signal) \\
rolling std division & volatility-normalizes across symbols and eras &
divides away volatility information --- keep a copy of the raw std as a
feature if vol matters \\
cross-sectional rank / z-score per timestamp & removes market-wide
component; robust to fat tails & discards the market component ---
reintroduce it explicitly as its own feature \\
quantile clip at fitted bounds & tames gradient-dominating outliers &
saturates genuine extremes (choose bounds wide: we use $[0.1\%, 99.9\%]$)\\
\bottomrule
\end{tabularx}
\end{center}

The composition we settled on, mirroring the 8th-place solution and audited
in Chapter~\ref{ch:leakage}: quantile clip $\to$ standardize (both fitted on
training dates, frozen); \emph{plus} engineered columns that reintroduce
deliberately removed information --- the per-symbol rolling deviation
(``feature minus its own trailing mean over 1000 steps'') keeps the fast
component while the slow level is dropped, and the cross-symbol market
average per timestamp reintroduces the commonality that per-symbol work
removes.

\begin{keybox}[: normalization as decomposition, not destruction]
The pattern behind every good transform above: \emph{split} a series into
components (level + deviation; market + idiosyncratic; scale + shape) and
hand the model each part it deserves, rather than \emph{deleting} a
component by normalizing it away. Whenever you write a normalizer, ask what
it removes and where that information will re-enter the pipeline --- if the
answer is ``nowhere,'' pause. Chapter~\ref{ch:smoothing} develops the same
principle for temporal filtering.
\end{keybox}

\section{Online statistics: the scaler is a model too}

Under drift, frozen statistics go stale: a feature whose variance doubles
mid-validation will be persistently mis-scaled by a frozen standardizer.
The deployable compromise, matching what a live system could do:
\textbf{freeze at the training boundary; then let running statistics update
from days already predicted} (Welford's algorithm makes this exact and
$O(1)$ per row). Our \code{OnlineStandardizer} does exactly this during the
walk-forward evaluation --- \code{partial\_fit} on day $d{-}1$ after day
$d{-}1$ is scored, never before.

One subtlety worth a sentence: if both your scaler \emph{and} your model
adapt online (Chapter~\ref{ch:online}), they co-drift. Refit-learning-rate
sweeps must therefore be run with the scaler behavior fixed to its
production setting, or the sweep answers a question about a system that
will not exist.

\section{Precision: float16 storage, float32 compute}

At 47M rows $\times$ 150 engineered columns, storage precision decides what
machines can hold the problem:

\begin{itemize}
\item float64: $\sim 56$ GB --- out of reach for laptops and free Colab.
\item float32: $\sim 28$ GB --- fits big workstations only.
\item \textbf{float16: $\sim 14$ GB} --- fits everywhere that matters.
\end{itemize}

The accuracy question: standardized features are $O(1)$ magnitudes, where
float16's relative error is $\sim 5 \times 10^{-4}$; round-tripping our
standardized matrix through float16 changed values by $\sim 10^{-4}$ ---
two orders of magnitude below feature noise. The rule that makes this safe:
\textbf{numerically sensitive operations (rolling statistics, market
averages, standardization) run in float32; only the finished, standardized
matrix is stored in float16; models upcast blocks to float32 at load}. We
measured no score difference between float16-stored and float32-stored
pipelines --- but the halved footprint converted several ``impossible''
experiments into overnight ones.

\section{Precompute once, memmap forever}

The single best workflow decision of the campaign. Feature engineering
needs the whole panel in memory (rolling windows cross day boundaries;
market averages need all symbols), takes minutes, and is identical across
every model. So: do it \emph{once}, and serve every experiment from the
result.

\begin{lstlisting}
# precompute_dataset.py (sketch of the layout it writes)
out/
  X.f16.mmap        # (N, K) float16 memmap - standardized features
  y.f32.npy         # target;  w.f32.npy  weights
  resp.f16.npy      # auxiliary-target matrix
  symbols/dates/times.i16.npy
  preprocessor.pkl  # frozen clipper + scaler (the leakage boundary)
  manifest.json     # feature names, per-date row ranges, config
\end{lstlisting}

Two design details carry most of the value:

\begin{itemize}
\item \textbf{Row order is (symbol, date, time).} Every (symbol, day) block
  is contiguous, so a sequence model materializes its
  (symbols $\times$ 968 $\times$ features) day-tensor with pure slicing ---
  no gather, no sort, no copy. Choosing row order to match the training
  access pattern is the difference between I/O-bound and compute-bound
  epochs.
\item \textbf{The manifest maps dates to row ranges.} Training on ``dates
  1319--1598'' is a dictionary lookup plus a memmap slice. This is what
  makes \emph{temporal bagging} (Chapter~\ref{ch:ensembling}) --- many
  models on many overlapping date subsets --- essentially free to
  orchestrate: every bag member reads its own date list from the same
  file, and peak RAM stays at the size of one member's slice.
\end{itemize}

The same precompute is the natural home for the leakage boundary: the
script takes \code{-{}-train-until}, fits the preprocessor strictly below
it, and transforms everything with frozen statistics
(Chapter~\ref{ch:leakage}). Downstream code \emph{cannot} re-leak what it
never recomputes.

\begin{jsbox}[: measured footprint]
Full pool, 999 usable dates, 34.8M rows $\times$ 134 columns: precompute
runs in 7--11 minutes at $\sim$5 GB peak RAM and writes a 9 GB float16
memmap. A training job on a 0.6-fraction date subsample then peaks around
9 GB --- inside a 16 GB laptop --- where the naive
``engineer-features-then-train'' loop needed $\sim$28 GB and half an hour
\emph{per experiment}. The memmap paid for itself on day one; every
number in Part~V was trained from it.
\end{jsbox}

\section{Fitting statistics without materializing the world}

One implementation trick generalizes: fitting the preprocessor does not
require the full matrix in memory. Quantile clip bounds are stable
statistics --- fit them on every $k$-th training day. Means and variances
stream exactly via Welford updates one day at a time. Our
\code{fit\_preprocessor\_streaming} holds one day of float32 at a time and
never allocates the pool; the pattern converts ``needs a big machine'' into
``needs any machine, patiently.''

\section{Checklist}

\begin{enumerate}
\item Profile pathologies first: raggedness, missing blocks, structured
  nulls, tails. Decide policies explicitly; write them down.
\item Choose normalizations per feature family by the invariance they buy;
  reintroduce removed components as explicit features.
\item Fit statistics on training dates only; freeze; online-update only
  from predicted-past days.
\item Store float16, compute float32; verify round-trip error against
  feature noise once.
\item Precompute the engineered matrix once; memmap; row-order by access
  pattern; date-range manifest for slicing.
\item Stream statistic fitting day-by-day; never hold the pool to fit a
  scaler.
\end{enumerate}
